% GPT-6 Astra Pro; /root/fixed_surface; 22 November 2026 UTC.
% Effort not exposed. Complete proof for independent review.
\section[HD-F1-S28: complete marked-divisor support and trivial point stabilizer]{HD-F1-S28: the marked-divisor support is exactly the covering curve, and the actual norm sheaf has trivial point stabilizer}
\label{sec:f1-norm-fourfold-marked-support-stabilizer}

Retain the actual smooth norm component $T$, its normal Abel
image $Y\subset P$, the map $a:T\to P$, and the constructions
of S13--S25, on their specified open with $C$ nonhyperelliptic.
The additional nontrigonal restriction in S26 can be removed.
All its marked-divisor cohomology and relative Hilbert-curve
conclusions hold on this open. Moreover, for every fixed $z\in Z$,
\[
 \begin{split}
 V^0(T,M_z)
  &:=\{\delta\in\operatorname{Pic}^0(P):
             H^0(T,M_z\otimes a^*\delta)\ne0\}\\
  &=\{\nu_z(z'):z'\in Z\},\qquad
          a^*\nu_z(z')=M_{z'}\otimes M_z^{-1}.
 \end{split}
 \tag{F1S28.1}
\]
Here $\nu_z$ is the actual closed immersion of the covering
curve specified below; the equality is an equality of closed
sets with their reduced structures. At every point of this
set the space of sections has dimension one. No assertion
about a nonreduced cohomology-jump scheme is needed.

The resulting full support calculation proves
\[
 \{x\in P:t_xY=Y\}=\{0\},\qquad
 \{(x,L)\in P\times\widehat P:
           t_x^*\mathcal O_Y\otimes L\simeq\mathcal O_Y\}
                      =\{(0,\mathcal O_P)\}.
 \tag{F1S28.2}
\]
Consequently every specified faithful finite point action
admitting a lift of this actual sheaf is the trivial group.
On the S23 period-rank-twelve open, S27's four-dimensional
trace-kernel conclusion also holds without the extra
nontrigonal condition.

\paragraph{1. The two paracanonical systems have no common base point.}
Write $\varepsilon$ for the nontrivial order-three line on
$C$ defining the original connected cover. Suppose a point
$p$ were a base point of both $K_C\varepsilon$ and
$K_C\varepsilon^{-1}$. Each system has four sections.
Riemann--Roch for $K_C\varepsilon(-p)$ gives
$h^0(\varepsilon^{-1}(p))=1$, and the inverse system gives
$h^0(\varepsilon(p))=1$. Thus actual effective degree-one
divisors $q,q'$ satisfy
\[
       \varepsilon=\mathcal O_C(p-q),\qquad
       \varepsilon^{-1}=\mathcal O_C(p-q'),\qquad
                              q+q'\sim2p.
 \tag{F1S28.3}
\]
On the nonhyperelliptic curve two distinct effective
degree-two divisors cannot be linearly equivalent: their
moving part would give a degree-at-most-two map to the
projective line. Therefore $q+q'=2p$ as divisors, so
$q=q'=p$ and $\varepsilon$ would be trivial, a contradiction.

It follows that at every $z\in Z$ at least one of the two
nontrivial eigenspaces of canonical forms evaluates
nontrivially. Those eigenspaces annihilate
$U=\pi^*H^1(\mathcal O_C)$ under Serre duality. Hence the
actual Abel tangent line $\ell_z$ satisfies
$\ell_z\cap U=0$, exactly the rank-five condition in
S26(F1S26.7)--(F1S26.9). This was the sole use of
nontrigonality in that proof. Its actual curve complexes,
Fourier calculation, component idempotents, excess triangle,
integral line identity and relative Hilbert construction
therefore give all of S26's conclusions here. In particular
$h^0(M_z)=1$, $h^1(M_z)=7$, and the actual section map in
degree one is surjective for every $z$.

\paragraph{2. The actual Picard parameter and the full support calculation.}
Fix an Abel map $j:Z\to JZ$ and the corresponding sum map
$u:X=\operatorname{Sym}^8Z\to JZ$. Retain
$F:X\to A=JC$ and $B=\operatorname{Pic}^0(A)$ from S26.
The composite $\operatorname{Nm}\,u$ differs from $F$ by
a fixed translation. The dual of the inclusion $P\hookrightarrow JZ$
is surjective. Thus every $\delta\in\operatorname{Pic}^0(P)$
has a lift $\gamma\in\operatorname{Pic}^0(JZ)$ with
\[
                              u^*\gamma|_T\simeq a^*\delta.
 \tag{F1S28.4}
\]
Set $\ell=j^*\gamma\in\operatorname{Pic}^0(Z)$.
These are actual pullbacks; no principal-polarization sign
is used to identify this parameter with a curve line.

Consider the actual curve complex on $B$
\[
 \mathcal C_{z,\gamma}=Rp_{B*}
  \bigl(\mathcal O_Z(z)\otimes\ell
                 \otimes(f\times1)^*\mathcal P\bigr),
 \tag{F1S28.5}
\]
where $\mathcal P$ and $f$ are the S26 normalized Poincare
line and curve map. A parameter of $B$ pulls back to
$\pi^*\alpha$ for an arbitrary $\alpha\in\operatorname{Pic}^0(C)$.
Suppose first that
\[
 H^0\bigl(Z,\mathcal O_Z(z)\otimes\ell\otimes\pi^*\alpha\bigr)=0
                   \quad\text{for every }\alpha.
 \tag{F1S28.6}
\]
The degree is one and $g(Z)=13$, so all the corresponding
$H^1$ spaces have dimension eleven. Cohomology and base
change for the proper smooth curve projection make
\[
                         \mathcal C_{z,\gamma}=\mathcal V[-1]
 \tag{F1S28.7}
\]
for a vector bundle $\mathcal V$ of rank eleven on $B$.
The exact Poincare addition, product Kunneth and permutation
idempotent maps of S26 now give
\[
 \Phi_A\bigl(RF_*(L_z\otimes u^*\gamma)\bigr)
          =\operatorname{Sym}^8_{\rm der}\mathcal C_{z,\gamma}
          =\bigwedge^8\mathcal V[-8].
 \tag{F1S28.8}
\]
Indeed pullback to $Z^8$ of the additional line $u^*\gamma$
is $(j^*\gamma)^{\boxtimes8}$, by the normalized addition
isomorphism for a degree-zero line. Thus the displayed
formula is the actual coefficient-twisted version of S26.

Mukai inversion in dimension five identifies the original
direct image with
\[
        (-1_A)^*\Phi_B(\bigwedge^8\mathcal V)[-3].
 \tag{F1S28.9}
\]
Every derived fibre of this object is concentrated in
degrees at least three. More explicitly, the integrand
on $B\times A$ is a vector bundle, hence flat over $A$;
proper derived base change gives the fibre
$R\Gamma(B,\bigwedge^8\mathcal V\otimes\mathcal P_a)[-3]$,
whose unshifted complex has no negative cohomology.

Apply this to the canonical point of $A$. The S21 excess
triangle for this derived fibre is the same triangle
twisted by the actual line $L_z\otimes u^*\gamma$.
Its excess line on each of the three norm components is
\[
 M_z\otimes a^*\delta\otimes\mathcal O_T(-h)
      =(M_{\sigma z}\otimes M_{\sigma^2z})^{-1}
                                      \otimes a^*\delta.
 \tag{F1S28.10}
\]
It is antiample: point-divisor lines are ample, and a
degree-zero twist does not change ampleness. Kodaira
vanishing and Serre duality give zero cohomology in
degrees below four. Thus degree zero of the derived
fibre is degree zero on the whole ordinary norm fibre.
By (F1S28.9) it is zero, so it is zero on each component.
This proves the required vanishing without selecting a
new component idempotent.

The contrapositive says that a nonzero section on the
specified component forces a section of one of the
degree-one curve lines in (F1S28.6). Consequently some
$z'\in Z$ and $\alpha\in\operatorname{Pic}^0(C)$ give
the actual equality
\[
                       j^*\gamma
                  =\mathcal O_Z(z'-z)\otimes\pi^*\alpha^{-1}.
 \tag{F1S28.11}
\]
The Abel pullback isomorphism
$j^*:\operatorname{Pic}^0(JZ)\to\operatorname{Pic}^0(Z)$
identifies its first factor with the line whose pullback
to $X$ is $L_{z'}L_z^{-1}$. To check the latter equality
integrally, pull back to $Z^8$: both become
$\mathcal O_Z(z'-z)^{\boxtimes8}$. The injectivity of
line pullback by $Z^8\to X$, proved in S26 using the
permutation action and the small diagonal, gives equality
on $X$ itself.

The second factor of (F1S28.11) is the pullback from
$JC$ along the norm. Its pullback to $X$ is therefore
$F^*$ of a degree-zero line, up to the fixed translation
which acts trivially on its Picard class. Since $F$ is
constant on $T$, that factor restricts trivially. Hence
\[
                         M_z\otimes a^*\delta\simeq M_{z'}.
 \tag{F1S28.12}
\]
Conversely every such line has its actual marked section,
and its section space has dimension one by paragraph1.
This proves (F1S28.1). The map $\nu_z$ is injective by
S26's degree-two divisor argument and has nonzero
differential by the projected Abel tangent computation.
Being proper, it is a closed immersion of the smooth
curve into $\operatorname{Pic}^0(P)$.

\paragraph{3. Actual translations preserve this Hilbert component.}
Let $H=\{x\in P:t_xY=Y\}$, the closed algebraic subgroup
obtained as the stabilizer of the actual Hilbert point
of $Y$. Every element has the unique regular lift
$f_x:T\to T$ from S20, with
\[
                a f_x=t_x a,\qquad f_x f_y=f_{x+y}.
 \tag{F1S28.13}
\]
One can equivalently use the intrinsic canonical
projectivization from S25. The S20 construction through
the actual pluricanonical ring already proves regularity,
uniqueness, and the group law for every stabilizing
translation.

The subgroup $H$ is finite. A positive-dimensional
identity component would contain a nonzero point of
order two. Its lifted involution would be free by
(F1S28.13). The resulting finite etale quotient of the
smooth fourfold would give
\[
              7371=\int_Tc_4(T_T)
                       =2\int_{T/\langle f_x\rangle}
                                      c_4(T_{T/\langle f_x\rangle}),
 \tag{F1S28.14}
\]
contradicting integrality. Thus $H$ is zero-dimensional,
and a closed subgroup of the projective abelian variety
with this dimension is finite. In particular each $f_x$
has finite order.

Fix $x\in H$, write $f=f_x$, and choose a representative
$\mathcal B$ of $\beta$. The actual line identity
S26(F1S26.18), together with $f^*\omega_T=\omega_T$
and (F1S28.13), gives, for every $z$,
\[
 f^*M_z=M_z\otimes a^*\delta_x,
 \qquad
 \delta_x=(t_x^*\mathcal B\otimes\mathcal B^{-1})^{-1}
                                      \in\operatorname{Pic}^0(P).
 \tag{F1S28.15}
\]
The degree-zero correction in that identity is unchanged
in Picard class by $t_x$. Also $f^*$ is identity on
$\operatorname{Pic}^0(T)$: the isomorphism
$a^*:\operatorname{Pic}^0(P)\to\operatorname{Pic}^0(T)$
intertwines it with translation pullback, which is
identity there.

For a fixed base marked point $z_0$, pullback of sections
and (F1S28.15) therefore make translation by $\delta_x$
preserve the whole support $V^0(T,M_{z_0})$. Paragraph2
identifies its reduced curve precisely, so it induces
an automorphism $\gamma$ of the actual $Z$. Equivalently,
the line $f^*M_z$ is a unique $M_{\gamma(z)}$. Since
both have one-dimensional section spaces, their effective
divisors agree:
\[
                              f^{-1}(D_z)=D_{\gamma(z)}.
 \tag{F1S28.16}
\]
This proves preservation of the original Hilbert component;
it does not presume that automorphisms preserve an arbitrary
component. The induced maps are morphisms because translation
on Picard space and the closed immersion $\nu_{z_0}$ are
morphisms. Taking powers in (F1S28.16) shows that the order
of $\gamma$ divides the finite order of $f$.

\paragraph{4. The induced curve action fixes the actual Prym pointwise.}
Let $\Gamma:JZ\to JZ$ be the Jacobian automorphism
induced by pushforward of divisors under $\gamma$.
The universal-divisor equality (F1S28.16) gives the exact
equivariance of the original inclusion $T\hookrightarrow
\operatorname{Sym}^8Z$:
\[
                         f(t)=\gamma^{-1}_*(t)
                       \quad\text{as effective divisors on }Z.
 \tag{F1S28.17}
\]
To include multiplicities, view (F1S28.16) as equality
of the two maps from $Z$ to its smooth Hilbert component.
Pulling back the universal family gives equality of the
relative Cartier divisors on $Z\times T$; viewing that
same family over $T$ gives (F1S28.17).

Normalize the Abel map on $T$ at a point $t_0$. Its
Albanese homomorphism is $U=iA$, where
$A:\operatorname{Alb}(T)\xrightarrow{\sim}P$ is the
actual S21 isomorphism and $i:P\hookrightarrow JZ$
is the original inclusion. Equation (F1S28.13) gives
$A f_{\rm Alb}=A$, hence $f_{\rm Alb}=1$. Taking the
Albanese homomorphism in (F1S28.17), with the constant
at $f(t_0)$ subtracted, gives
\[
                    U f_{\rm Alb}=\Gamma^{-1}U,
                         \qquad \Gamma i=i.
 \tag{F1S28.18}
\]
Thus it is the literal embedded Prym that is fixed
pointwise.

Suppose $\gamma\ne1$, and let $m\ge2$ be its finite
order. The quotient map $\rho:Z\to B_0=Z/\langle\gamma\rangle$
is a finite morphism of smooth projective curves of
degree $m$. The pullback of invariant differentials
identifies
\[
           H^0(B_0,K_{B_0})
                  =H^0(Z,K_Z)^{\langle\gamma\rangle}.
 \tag{F1S28.19}
\]
At a ramification point this follows directly from
the coordinates $t=z^e$: an invariant regular form
has the shape $z^{e-1}g(z^e)\,dz$ and descends to
$g(t)\,dt/e$. Thus this equality also retains all
ramification points.

The invariant space in (F1S28.19) has dimension at
least eight. Indeed $i^*:H^0(JZ,\Omega^1)\to
H^0(P,\Omega^1)$ is surjective; averaging any lift
of a form on $P$ under $\Gamma$ leaves its restriction
unchanged by (F1S28.18). The Abel pullback identifies
these invariant forms with those on $Z$. Hence
$g(B_0)\ge8$. Riemann--Hurwitz now gives the impossible
inequality
\[
          24=2g(Z)-2=m(2g(B_0)-2)+\deg R
                                      \ge2(16-2)=28.
 \tag{F1S28.20}
\]
Therefore $\gamma=1$. Equation (F1S28.17) and the
actual closed immersion into the symmetric product
give $f=1$. Equation (F1S28.13), evaluated at any
point of $T$, then gives $x=0$. This proves the first
equality in (F1S28.2).

Finally an isomorphism
$t_x^*\mathcal O_Y\otimes L\simeq\mathcal O_Y$
forces the same support, hence $x=0$. It then gives
$L|_Y\simeq\mathcal O_Y$, so $a^*L$ is trivial.
The actual S19 Picard injection forces
$L=\mathcal O_P$. This proves the second equality
and the assertion about every faithful finite point
action. Its coherence data are unchanged; faithfulness
alone makes a group with this singleton point image
trivial.

\paragraph{Exact source and deformation scope.}
The Fourier transforms, normalized Poincare addition,
inverse shift and flat derived-fibre calculation are
the actual maps of S26 and the same checked Mukai
Proposition1.3 and Theorem2.2. The only additional
cohomology computation is the global vector bundle
in (F1S28.7); its inverse transform is evaluated
explicitly in (F1S28.9). The norm-factor pullback in
(F1S28.11)--(F1S28.12) specifies the required Picard
comparison without a polarization-sign convention.
The lifted automorphism and its free quotient are
the already proved S20 maps. The final quotient-curve
calculation writes both invariant differential descent
and the complete ramification term.

All other assumptions of S13--S25 are retained. This
does not extend their smoothness or pencil hypotheses
to a boundary where they were not proved. On their
intersection with the S23 period-rank open, paragraph1
supplies S26 to the actual reconstruction in S27; its
rank-four obstruction and full trace-kernel statements
therefore hold there without the extra nontrigonal
restriction. No equality between the full self-Ext
algebra and an exterior quotient is asserted.
